\documentclass[10pt,a4paper]{article}
\usepackage{amsmath,amssymb,amsthm,graphicx}
\title{Characteristic-free constructive set-multilinearization preserving low product depth}
\author{Michael A. Forbes}
\date{doi:10.4230/LIPIcs.CCC.2024.31; CC BY 4.0}
\begin{document}
\maketitle
\noindent\textit{Allowed standard prerequisites (editorial).}\begin{itemize}
\item Binet-Minc identity for rectangular permanents. Explicitly quoted author Theorem13 citing Minc[17]; full identity retained.
\item Bell-number asymptotics and elementary partition counting. Author Fact15 citing Flajolet,Sedgewick[6], Definition14 and Lemma16; exact statements retained.
\end{itemize}\medskip
\section*{Author circuit and polynomial conventions}
Algebraic complexity asks for the best methods to compute multivariate polynomials $f(x_1,\ldots,x_n)$ from the primitive operations of addition and multiplication, using constants from the underlying field $\mathbb F$ for free. These computations are organized as an algebraic circuit, a directed acyclic graph of these primitive operations from the input variables $x_1,\ldots,x_n$. The complexity of this object is primarily governed by its size, which is the number of operations in the circuit. Another complexity measure is the product depth, which is the maximum number of product operations on any input to output path.
The iterated matrix multiplication polynomial $\operatorname{IMM}_{n,d}=(X_1\cdots X_d)_{1,1}$ where $X_i$ is an $n\times n$ symbolic matrix. $F_k$ is the $k$-th Fibonacci number (so $F_0=0$, $F_1=F_2=1$, $F_3=2$, etc.).

\paragraph{Definition 7.} Let the variables $\bar x$ be partitioned into $\bar x=\bar x_1\sqcup\cdots\sqcup\bar x_d$. A monomial is set-multilinear (with respect to the partition $\bar x=\bar x_1\sqcup\cdots\sqcup\bar x_d$) if it can be written has $\prod_{i=1}^d(\bar x_i)_{j_i}$ for some $j_1,\ldots,j_d$.
Define the set-multilinear projection (with respect to the partition $\bar x=\bar x_1\sqcup\cdots\sqcup\bar x_d$) to be the linear map $\pi_{\rm sm}:\mathbb F[\bar x]\to\mathbb F[\bar x]$ which is identity on set-multilinear monomials, and zero on all other monomials.

\paragraph{Definition 11.} Let $X$ be an $n\times m$ symbolic matrix of variables, with $n\le m$, where $X_{i,j}=x_{i,j}$ are distinct variables. Define the (rectangular) permanent $\operatorname{perm}(X)\in\mathbb Z[(x_{i,j})_{i\in[n],j\in[m]}]$ by
\[\operatorname{perm}_{n\times m}(X)=\sum_{\sigma:[n]\hookrightarrow[m]}x_{1,\sigma(1)}\cdots x_{n,\sigma(n)},\]
that is, the sum runs over all injective maps $\sigma$ from $[n]$ to $[m]$.
Note in particular that $\operatorname{perm}_{n\times m}$ is a set-multilinear polynomial when we partition the matrix into its rows.
It is sometimes helpful to view the rectangular permanent as a sum of square permanents.
\paragraph{Lemma 12.} Let $X$ be an $n\times m$ symbolic matrix of variables, with $n\le m$, where $X_{i,j}=x_{i,j}$ are distinct variables. Then,
\[\operatorname{perm}_{n\times m}(X)=\sum_{S\in\binom{[m]}n}\operatorname{perm}_{n\times n}(X|_{[n]\times S}).\]
The above lemma shows the rectangular permanent is computable by $\operatorname{poly}(m^n)$ size algebraic circuits. However, the following identity gives a non-trivially better algorithm.
\paragraph{Theorem 13 (Binet-Minc Identity [17]).} Let $X$ be an $n\times m$ symbolic matrix of variables, with $n\le m$, then the permanent can be computed by
\[\operatorname{perm}_{n\times m}(X)=\sum_{\mathcal F\in\mathcal P_n}(-1)^{n-|\mathcal F|}\prod_{S\in\mathcal F}(|S|-1)!\sum_{j=1}^m\prod_{i\in S}x_{i,j}\]
where $\mathcal P_n$ is the set of all partitions of $[n]$ into (non-empty) sets, and $|\mathcal F|$ is the number of parts in the partition $\mathcal F$ of $[n]$.
To understand the complexity of this expression it is helpful to have the following definition.
\paragraph{Definition 14.} Define the $n$-th Bell number $B_n$ to be the number of ways to partition $[n]$ into (non-empty) sets.
We will use the following asymptotic estimate of Bell number size.
\paragraph{Fact 15 ([6]).} $B_n=\Theta(\frac n{\ln n})^n$.
The following lemma is easy to prove from the definition of Bell numbers.
\paragraph{Lemma 16.} $B_n\cdot B_m\le B_{n+m}$.
The Binet-Minc identity immediately implies the following algebraic circuit for the permanent.
\paragraph{Corollary 17.} The rectangular permanent $\operatorname{perm}_{n\times m}$ has a $\Sigma^{B_n}\Pi^n\Sigma^m\Pi^n$ formula of size $\operatorname{poly}(m,B_n)$, where the super-scripts are upper bounds on the respective fan-ins of the formula. Further, this formula is set-multilinear (and hence homogeneous) with respect to the partition of the $n\times m$ variables into rows.
\paragraph{Definition 20.} Let $X$ be an $n\times m$ symbolic matrix of variables, with $n\le m$, where $X_{i,j}=x_{i,j}$ are distinct variables. Let $k\le n$. Define the $k$-surjective (rectangular) permanent $\operatorname{perm}_{n\times m;k}(X)\in\mathbb Z[(x_{i,j})_{i\in[n],j\in[m]}]$ by
\[\operatorname{perm}_{n\times m;k}(X)=\sum_{\sigma:[n]\hookrightarrow[m],\operatorname{im}(\sigma)\supseteq[k]}x_{1,\sigma(1)}\cdots x_{n,\sigma(n)},\]
that is, the sum runs over all injective maps $\sigma$ from $[n]$ to $[m]$ that contain $[k]$ in their image.
Note in particular that $\operatorname{perm}_{n\times m;k}$ is a set-multilinear polynomial when we partition the matrix into its rows.
We can compute the surjective permanent by standard permanents.
\paragraph{Definition 23.} Let the variables $\bar x$ be partitioned into $\bar x=\bar x_1\sqcup\cdots\sqcup\bar x_d$.
Let $S\subseteq[d]$. A monomial is $S$-set-multilinear if it can be written as $\prod_{i\in S}(\bar x_i)_{j_i}$ for some $(j_i)_{i\in S}$. Define the $S$-set-multilinear projection to be the linear map $\pi_{{\rm sm},S}$ which is identity on $S$-set-multilinear monomials, and zero on all other monomials.
The set of $S$-set-multilinear monomials for some $S\subseteq[d]$ are called at most set-multilinear monomials. Define the non-set-multilinear projection to be the linear map $\pi_{\neg{\rm sm}}$ which is zero on at-most set-multilinear monomials, and identity on all other monomials.
In particular, for $S=[d]$, $\pi_{\rm sm}(f)$ and $\pi_{{\rm sm},S}(f)$ are the same. For $S=\varnothing$, $\pi_{{\rm sm},S}(f)$ is the constant part of $f$, $\pi_{{\rm sm},\varnothing}(f)=f(0)$. More generally, we can decompose $f$ into its set-multilinear parts as follows.

\paragraph{Corollary 27.} Let $\mathbb F$ be any field. Let the variables $\bar x$ be partitioned into $\bar x=\bar x_1\sqcup\cdots\sqcup\bar x_d$. Suppose $f\in\mathbb F[\bar x]$ can be computed by a size $s$ product-depth $\Delta$ algebraic circuit. Then the set-multilinear projection $\pi_{\rm sm}(f)\in\mathbb F[\bar x]$ can be computed by a size $\operatorname{poly}(s,\Theta(\frac d{\ln d})^d)$-size product-depth $2\Delta$ circuit.

\section*{Complete author-local construction supporting Corollary 27}
\paragraph{Lemma 21.} Let $X$ be an $n\times m$ symbolic matrix of variables, with $n\le m$, where $X_{i,j}=x_{i,j}$ are distinct variables. Let $k\le n$. Then, the $k$-surjective permanent can be written as
\[\operatorname{perm}_{n\times m;k}(X)=\sum_{S\in\binom{[n]}k}\operatorname{perm}_{k\times k}(X|_{S\times[k]})\cdot\operatorname{perm}_{(n-k)\times(m-k)}(X|_{([n]\setminus S)\times([m]\setminus[k])}).\]
Further, this expression is set-multilinear with respect to partitioning $X$ by its rows.
\begin{proof}
This follows from the observation that each map $\sigma:[n]\hookrightarrow[m]$ with $\operatorname{im}(\sigma)\supseteq[k]$ uniquely decomposes into a bijection $\tau:S\leftrightarrow[k]$ for some $S\in\binom{[n]}k$, and an injection $\nu:([n]\setminus S)\hookrightarrow([m]\setminus[k])$.
The claim about set-multilinearity then follows from noting that the multiplication $\operatorname{perm}_{k\times k}(X|_{S\times[d]})\cdot\operatorname{perm}_{(n-k)\times m-k}(X|_{([n]\setminus S)\times([m]\setminus[k])})$ is a product of two polynomials who only use disjoint rows of $X$.
\end{proof}
We now analyze the complexity of computing the surjective permanent, by reduction to standard permanents, and then applying the Binet-Minc identity.
\paragraph{Corollary 22.} Let $k\le n\le m$. The $n\times m$ $k$-surjective permanent has a $\operatorname{poly}(m,\Theta(\frac n{\ln n})^n)$-size depth-4 formula that is set-multilinear with respect to rows.
\begin{proof}
Via the above lemma, and the formula complexity of the Binet-Minc identity,
\begin{align*}
\operatorname{perm}_{n\times m;k}(X)&=\sum_{S\in\binom{[n]}k}\underbrace{\operatorname{perm}_{k\times k}(X|_{S\times[k]})}_{\Sigma^{B_k}\Pi^k\Sigma^k\Pi^k}\cdot\underbrace{\operatorname{perm}_{(n-k)\times(m-k)}(X|_{([n]\setminus S)\times([m]\setminus[k])})}_{\Sigma^{B_{n-k}}\Pi^{n-k}\Sigma^{m-k}\Pi^{n-k}},
\end{align*}
distributing the multiplication past the addition,
\[=\underbrace{\sum^{\binom nk}\sum^{B_k\cdot B_{n-k}}}_{\Sigma^{\binom nk B_k B_{n-k}}}\underbrace{\left(\prod^k\sum^k\prod^k\right)\cdot\left(\prod^{n-k}\sum^{m-k}\prod^{n-k}\right)}_{\Pi^n\Sigma^m\Pi^n}\]
from which the size bound follows by noting that $\binom nk B_k B_{n-k}\le2^n B_n\le\Theta(\frac n{\ln n})^n$ via Lemma 16 and Fact 15.
The set-multilinearity of this formula follows from the fact that the decomposition used here from the above lemma is set-multilinear, that Binet-Minc is set-multilinear, and that our use of the distributive law preserves set-multilinearity.
\end{proof}
We now proceed to give a non-trivial set-multilinearization for low-depth algebraic circuits. To do so, it will be helpful to have more notation for extracting various set-multilinear components of polynomials.
\paragraph{Lemma 24.} Let the variables $\bar x$ be partitioned into $\bar x=\bar x_1\sqcup\cdots\sqcup\bar x_d$. Then for any polynomial $f$,
\[f=\pi_{\neg{\rm sm}}(f)+f(0)+\sum_{\varnothing\ne S\subseteq[d]}\pi_{{\rm sm},S}(f).\]
It immediately follows that we can simulate an addition (or even a linear combination) of polynomials by instead adding the constituent set-multilinear parts.
\paragraph{Lemma 25.} Let $f$ be any field. Let the variables $\bar x$ be partitioned into $\bar x=\bar x_1\sqcup\cdots\sqcup\bar x_d$. Let $f_1,\ldots,f_m$ be polynomials, with $f=\alpha_1f_1+\cdots+\alpha_m f_m$ for $\alpha_1,\ldots,\alpha_m\in\mathbb F$. Let $S\subseteq[d]$. Then $\pi_{{\rm sm},S}(f)$ can be computed by a depth-1 $\operatorname{poly}(m,2^d)$-size set-multilinear $\Sigma$-circuit given $\{\pi_{{\rm sm},S}(f_j)\}_{S\subseteq[d],j\in[m]}$ as inputs.
Less trivially is the simulation of a multiplication, for which we use the formula for the surjective permanent (Corollary 22).
\paragraph{Lemma 26.} Let the variables $\bar x$ be partitioned into $\bar x=\bar x_1\sqcup\cdots\sqcup\bar x_d$. Let $f_1,\ldots,f_m$ be polynomials, with $f=f_1\cdot\cdots\cdot f_m$. Let $S\subseteq[d]$ be of size $\ell$. Then $\pi_{{\rm sm},S}(f)$ can be computed by a depth-4 $\operatorname{poly}(m,\Theta(\frac\ell{\ln\ell})^\ell,2^d)$-size set-multilinear $\Sigma\Pi\Sigma\Pi$-circuit given $\{\pi_{{\rm sm},S}(f_j)\}_{S\subseteq[d],j\in[m]}$ as inputs.
\begin{proof}
The number of inputs to the circuit is $m2^d$. It remains to bound the number of gates by $\operatorname{poly}(m,\Theta(\frac\ell{\ln\ell})^\ell)$.
Rearrange the $f_i$ as needed so that $f_1(0)=\cdots=f_k(0)=0$ and $f_{k+1}(0),\ldots,f_m(0)\ne0$, for some $k\le m$. Thus, we can normalize the computation via
\[f=\prod_{i>k}^m f_i(0)\cdot\prod_{i\in[k]}f_i\cdot\prod_{i>k}^m\frac{f_i}{f_i(0)},\]
and thus define
\[g_i=\begin{cases}f_i&i\le k\\f_i/f_i(0)&i>k\end{cases},\]
and $g=\prod_i g_i$. It then follows that $f=\prod_{i>k}^m f_i(0)\cdot g$, $g_{k+1}(0)=\cdots=g_m(0)=1$, and $\pi_{{\rm sm},S}(f)=\prod_{i>k}^m f_i(0)\cdot\pi_{{\rm sm},S}(g)$. As $\prod_{i>k}^m f_i(0)$ is a non-zero constant, it suffices to prove the claim for $\pi_{{\rm sm},S}(g)$.
Now write $g$ as
\[g=\prod_{i=1}^k\left(\pi_{\neg{\rm sm}}(g_i)+\sum_{\varnothing\ne S_i\subseteq[d]}\pi_{{\rm sm},S_i}(g_i)\right)\cdot\prod_{i>k}^m\left(\pi_{\neg{\rm sm}}(g_i)+1+\sum_{\varnothing\ne S_i\subseteq[d]}\pi_{{\rm sm},S_i}(g_i)\right).\]
In expanding the above product, the only at-most set-multilinear terms are of the form
\[\prod_j\pi_{{\rm sm},S_{i_j}}(g_{i_j}),\]
where the $S_{i_j}\subseteq[d]$ are disjoint (as otherwise we create non-multilinear terms), and the indices $\{i_j\}_j$ are distinct (as we take exactly one term from each $g_i$ [possibly the term 1]).
Further, as $g_i(0)=0$ for $i\le k$, we must have $\{i_j\}_j\supseteq[k]$. Hence, by collecting like terms, we see that
\[\pi_{{\rm sm},S}(g)=\sum_{\mathcal F\in\mathcal P_S}\sum_{\sigma:\mathcal F\hookrightarrow[m];\operatorname{im}(\sigma)\supseteq[k]}\prod_{T\in\mathcal F}\pi_{{\rm sm},T}(g_{\sigma(T)})\]
where $\mathcal P_S$ is the collection of set partitions of $S$, which we can rewrite in terms of $k$-surjective permanents as
\[=\sum_{\mathcal F\in\mathcal P_S}\operatorname{perm}_{|\mathcal F|\times m;k}(A_{\mathcal F}),\]
where for a partition $\mathcal F$ of $S$ the matrix $A_{\mathcal F}$ is $|\mathcal F|\times m$ size matrix defined by
\[(A_{\mathcal F})_{T,j}=\pi_{{\rm sm},T}(g_j),\]
for $T\in\mathcal F$ and $j\in m$. Note that the matrix $A_{\mathcal F}$ has rows that access disjoint parts of the set-multilinear partition of $\bar x$, and so as the $k$-surjective permanent is set-multilinear with respect to rows, the computation $\operatorname{perm}_{|\mathcal F|\times m;k}(A_{\mathcal F})$ is set-multilinear with respect to the partition of $\bar x$. Further, a set-multilinear computation of this permanent will be set-multilinear with respect to $\bar x$. Hence, we can use the set-multilinear computation from Corollary 22 for the surjective permanent, which in this case yields a depth-4 set-multilinear formula of size $\operatorname{poly}(m,\Theta(\frac{|\mathcal F|}{\ln|\mathcal F|})^{|\mathcal F|})$.
Computing $\pi_{{\rm sm},S}(g)$ is then a sum of $B_{|S|}$ many such $\operatorname{perm}_{|\mathcal F|\times m;k}(A_{\mathcal F})$ terms, which does not increase depth as we collapse two sequential layers of addition gates, and this preserves set-multilinearity. The resulting size of the expression is $\operatorname{poly}(m,\Theta(\frac{|\mathcal F|}{\ln|\mathcal F|})^{|\mathcal F|},B_{|S|},2^d))\le\operatorname{poly}(m,\Theta(\frac\ell{\ln\ell})^\ell,2^d)$, as $|S|=\ell$, $|\mathcal F|\le|S|$.
\end{proof}
We now conclude with our set-multilinearization result over any field by gate-simulation.
\begin{proof}[Proof of Corollary 27]
By gate simulation. Let $\Phi$ be the hypothesized circuit computing $f$. For each node $v$ in $\Phi$, split $v$ into its at-most set-multilinear parts $\pi_{{\rm sm},S}(v)$ for each $S\subseteq[d]$.
If $v=\alpha_1v_1+\cdots+\alpha_m v_m$, then we can express the at-most set-multilinear parts of $v$ in terms of the $v_i$ using a $\operatorname{poly}(m,2^d)$-size product-depth 0 set-multilinear circuit by Lemma 25. If $v=v_1\times\cdots\times v_m$, then we can express the at-most set-multilinear parts of $v$ in terms of the $v_i$ using a $\operatorname{poly}(m,\Theta(\frac d{\ln d})^d)$-size product-depth 2 set-multilinear circuit by Lemma 26.
Correctness of the computation follows by induction on the circuit.
The overall size of the circuit has increased from $s$ to $\operatorname{poly}(s,\Theta(\frac d{\ln d})^d)$ by counting the size of the additional local gadgets of the gate simulation. Simulation of addition gates adds no product depth. Each product gate in the original circuit is turned into a product-depth 2 circuit in the gate simulation, hence the overall product-depth has at most doubled.
\end{proof}

\begin{thebibliography}{99}
\bibitem{source1} Prashanth Amireddy, Ankit Garg, Neeraj Kayal, Chandan Saha, and Bhargav Thankey. Low- depth arithmetic circuit lower bounds: Bypassing set-multilinearization. In Proceedings of the 50th International Colloquium on Automata, Languages and Programming (ICALP 2023), volume 261 of Leibniz International Proceedings in Informatics (LIPIcs), pages 12:1–12:20, 2023. Full version in the Electronic Colloquium on Computational Complexity (ECCC), Technical Report TR22-151. doi:10.4230/LIPICS.ICALP.2023.12.
\bibitem{source2} Robert Andrews. Algebraic hardness versus randomness in low characteristic. In Proceedings of the 35th Annual Computational Complexity Conference (CCC 2020), volume 169 of Leibniz International Proceedings in Informatics (LIPIcs), pages 37:1–37:32, 2020. Full version at arXiv:2005.10885. doi:10.4230/LIPICS.CCC.2020.37.
\bibitem{source3} C. S. Bhargav, Sagnik Dutta, and Nitin Saxena. Improved lower bound, and proof barrier, for constant depth algebraic circuits. In Proceedings of the 47th Internationl Symposium on the Mathematical Foundations of Computer Science (MFCS 2022), volume 241 of Leibniz International Proceedings in Informatics (LIPIcs), pages 18:1–18:16, 2022. doi:10.4230/ LIPICS.MFCS.2022.18.
\bibitem{source4} Radu Curticapean, Nutan Limaye, and Srikanth Srinivasan. On the VNP-hardness of some monomial symmetric polynomials. In 42nd International Conference on Foundations of Software Technology and Theoretical Computer Science (FSTTCS 2022), volume 250 of Leibniz International Proceedings in Informatics (LIPIcs), pages 16:1–16:14, 2022. Full version in the Electronic Colloquium on Computational Complexity (ECCC), Technical Report TR22-139. doi:10.4230/LIPICS.FSTTCS.2022.16.
\bibitem{source5} Ismor Fischer. Sums of like powers of multivariate linear forms. Mathematics Magazine, 67(1):59–61, 1994. URL: http://www.jstor.org/stable/2690560.
\bibitem{source6} Philippe Flajolet and Robert Sedgewick. Analytic combinatorics. Cambridge University Press, 2009. doi:10.1017/CBO9780511801655.
\bibitem{source7} Michael A. Forbes, Amir Shpilka, and Ben Lee Volk. Succinct hitting sets and barriers to proving algebraic circuits lower bounds. In Proceedings of the 49th Annual ACM Symposium on Theory of Computing (STOC 2017), pages 653–664, 2017. Full version at arXiv:1701.05328. doi:10.1145/3055399.3055496.
\bibitem{source8} Hervé Fournier, Nutan Limaye, Srikanth Srinivasan, and Sébastien Tavenas. On the power of homogeneous algebraic formulas. Electronic Colloquium on Computational Complexity (ECCC), TR23-191, 2023. URL: https://eccc.weizmann.ac.il/report/2023/191.
\bibitem{source9} Nashlen Govindasamy, Tuomas Hakoniemi, and Iddo Tzameret. Simple hard instances for low-depth algebraic proofs. In Preliminary version in the 63rd Annual IEEE Symposium on Foundations of Computer Science (FOCS 2022), pages 188–199, 2022. Full version at arXiv:2205.07175. doi:10.1109/FOCS54457.2022.00025.
\bibitem{source10} Joshua A. Grochow, Mrinal Kumar, Michael E. Saks, and Shubhangi Saraf. Towards an algebraic natural proofs barrier via polynomial identity testing. arXiv, 1701.01717, 2017. URL: http://arxiv.org/abs/1701.01717.
\bibitem{source11} Joshua A. Grochow and Toniann Pitassi. Circuit complexity, proof complexity, and polynomial identity testing. In Proceedings of the 55th Annual IEEE Symposium on Foundations of Computer Science (FOCS 2014), pages 110–119, 2014. Full version at arXiv:1404.3820. doi:10.1109/FOCS.2014.20.
\bibitem{source12} Alan Guo, Swastik Kopparty, and Madhu Sudan. New affine-invariant codes from lifting. In Proceedings of Innovations in Theoretical Computer Science (ITCS 2013), pages 529–540, 2013. Full version at arXiv:1208.5413. doi:10.1145/2422436.2422494.
\bibitem{source13} Ankit Gupta, Pritish Kamath, Neeraj Kayal, and Ramprasad Saptharishi. Arithmetic circuits: A chasm at depth three. In Proceedings of the 54th Annual IEEE Symposium on Foundations of Computer Science (FOCS 2013), pages 578–587, 2013. Full version in the Electronic Colloquium on Computational Complexity (ECCC), Technical Report TR13-026. doi:10. 1109/FOCS.2013.68.
\bibitem{source14} Erich L. Kaltofen. Factorization of polynomials given by straight-line programs. In Silvio Micali, editor, Randomness and Computation, volume 5 of Advances in Computing Research, pages 375–412. JAI Press, Inc., Greenwich, CT, USA, 1989. URL: http://www.math.ncsu. edu/~kaltofen/bibliography/89/Ka89\_slpfac.pdf.
\bibitem{source15} Nutan Limaye, Srikanth Srinivasan, and Sébastien Tavenas. Superpolynomial lower bounds against low-depth algebraic circuits. In Preliminary version in the 62nd Annual IEEE Sym- posium on Foundations of Computer Science (FOCS 2021), pages 804–814, 2022. Full version in the Electronic Colloquium on Computational Complexity (ECCC), Technical Report TR21-081. doi:10.1109/FOCS52979.2021.00083.
\bibitem{source16} Nutan Limaye, Srikanth Srinivasan, and Sébastien Tavenas. Guest column: Lower bounds against constant-depth algebraic circuits. SIGACT News, 53(2):40–62, 2022. doi:10.1145/ 3544979.3544989.
\bibitem{source17} Henryk Minc. Evaluation of permanents. Proc. Edinburgh Math. Soc. (2), 22(1):27–32, 1979. doi:10.1017/S0013091500027760.
\bibitem{source18} Ran Raz. Elusive functions and lower bounds for arithmetic circuits. Theory of Computing, 6(1):135–177, 2010. 40th Annual ACM Symposium on Theory of Computing (STOC 2008). doi:10.4086/TOC.2010.V006A007.
\bibitem{source19} Alexander A. Razborov. Lower bounds on the size of bounded depth circuits over a complete basis with logical addition. Matematicheskie Zametki, 41(4):598–607, April 1987.
\bibitem{source20} Amir Shpilka and Avi Wigderson. Depth-3 arithmetic circuits over fields of characteristic zero. Computational Complexity, 10(1):1–27, 2001. Preliminary version in the 14th Annual IEEE Conference on Computational Complexity (CCC 1999). doi:10.1007/PL00001609.
\bibitem{source21} Roman Smolensky. Algebraic methods in the theory of lower bounds for boolean circuit complexity. In Proceedings of the 19th Annual ACM Symposium on Theory of Computing (STOC 1987), pages 77–82, 1987. doi:10.1145/28395.28404.
\bibitem{source22} Ryan Williams. Finding paths of length k in O∗ (2k ) time. Inf. Process. Lett., 109(6):315–318, 2009. doi:10.1016/J.IPL.2008.11.004.
\end{thebibliography}
\end{document}
